\appendix

\section{Deployment Sites}
\label{app:sites}

% Recommendation: rewrite Sites as setting description and restate the
% load-bearing facts about each place in our own voice: what it looks like,
% who is there, who maintains it, its reputation, how we got access, and
% where the team stood.
\textbf{Tappen Park (Staten Island)} is a neighborhood park in Stapleton near a commercial district and the waterfront. Visitors use the park largely for sitting, talking, and spending time rather than organized recreation. Fixed trash receptacles are located around the park but are not distributed continuously through the open seating areas. The surrounding area also carries a documented reputation for public-safety concerns, including drug activity~\cite{sifentanyltaskforce2024}. We deployed TrashBot around the open lawn and paved paths, with the operator seated on a nearby bench.

\textbf{Astoria Park (Queens)} is a large riverside leisure park on the East River used for running, cycling, family outings, sunbathing, and other outdoor recreation. We deployed TrashBot near the swimming-pool entrance, along the riverside walk, and beside the skatepark. These areas placed the robot among parents with children, cyclists, and people sitting or lying on lawns and benches. Fixed trash receptacles were distributed across the park but often required walking some distance from the recreation areas. The site was covered by the same Parks Department permit, and the operator worked from a bench at the edge of the deployment area.

\textbf{Little Italy (Bronx)} is a dense commercial corridor centered on Arthur Avenue in Belmont and historically associated with the borough's Italian immigrant community~\cite{apa2016arthuravenue, helmreich2023bronx}. Restaurants, shops, market activity, and heavy pedestrian traffic occupy the corridor, and sanitation is visibly organized through local cleaning crews, porters, and fixed street receptacles. We deployed TrashBot in the plaza outside a pizzeria on the corridor. Access was arranged through the executive director of the neighborhood business association, who introduced the research team to nearby business owners.

\textbf{Campus Café (Manhattan)}The Manhattan deployment took place at a café beside a research institute. The café and its outdoor seating area are used by institute affiliates, local residents, and visitors. Deployment permission were granted by the institute's campus safety department.

\textbf{Brooklyn Navy Yard (Brooklyn)} is an access-controlled industrial and commercial campus used primarily by people working at organizations based there. Access was arranged
through a member of a community engagement group, who initially invited the research team to a community event. Permission to deploy TrashBot at the site was subsequently granted.


